\documentclass[11pt]{article}
\usepackage[T1]{fontenc}
\usepackage{lmodern,microtype,amsmath,amssymb,amsthm,mathtools,booktabs,array}
\usepackage[margin=1in]{geometry}
\usepackage{xcolor}
\usepackage[colorlinks=true,linkcolor=blue!45!black,citecolor=blue!45!black,urlcolor=blue!45!black]{hyperref}
\usepackage[numbers,sort&compress]{natbib}
\usepackage{enumitem}
\setlist{itemsep=3pt,topsep=5pt}
\setlength{\emergencystretch}{2em}
\newtheorem{theorem}{Theorem}[section]
\newtheorem{lemma}[theorem]{Lemma}
\newtheorem{proposition}[theorem]{Proposition}
\newtheorem{corollary}[theorem]{Corollary}
\theoremstyle{remark}\newtheorem{remark}[theorem]{Remark}
\newcommand{\TV}{\operatorname{TV}}
\newcommand{\E}{\mathbb E}
\newcommand{\Pbb}{\mathbb P}
\newcommand{\R}{\mathbb R}
\newcommand{\cP}{\mathcal P}
\newcommand{\1}{\mathbf 1}
\newcommand{\eps}{\varepsilon}
\newcommand{\KL}{\operatorname{KL}}
\DeclareMathOperator{\supp}{supp}
\title{Next-Symbol Prediction from Adversarially Corrupted Markov Trajectories}
\author{}
\date{2 November 2026}
\begin{document}
\maketitle
\begin{abstract}
A predictor observes one trajectory of an unknown homogeneous second-order finite-alphabet source after an adversary has replaced at most $k$ symbols. It must predict the source's next-symbol law at the true terminal context, not at an independent query. For fixed alphabet size $K$, history length $M+2$, and an all-state two-step minorization coefficient $\beta$, we determine the minimax total-variation risk when the final pair is protected. The rate is the larger of $M^{-1/2}$ and $(k/M)[1+\log(1+1/(ke_K))]$, capped at one, where $e_K=\max\{\beta,M^{-1}\}$ for $K\ge3$ and $e_2=\max\{\sqrt\beta,M^{-1}\}$. It holds for stationary and arbitrary initialization. The lower bounds account for the entire corrupted history, using common source-independent strict-budget kernels; the binary argument uses entropy transport to canonical Markov bridges. If the terminal symbols are editable but the target remains the true hidden terminal law, the rate instead contains the additional term $(1-\beta)\1_{\{k\ge1\}}$. The results are statistical full-information statements. Foundational clean prediction and specified occupation inequalities are inherited explicitly; no neural implementation or learning theorem is asserted.
\end{abstract}
\tableofcontents
\newpage
\section{Introduction}
Predicting the next symbol of a dependent process differs from estimating its transition table uniformly. The query is the terminal context of the same trajectory used for estimation. A context that appears only near the end can be statistically important even though it has few earlier departures. Adversarial replacement of recorded symbols adds a second difficulty: a changed symbol alters several overlapping triples and can change which observations appear to belong to the queried row.

This paper studies these effects without an intermediate information bottleneck. The predictor receives every corrupted symbol and may perform arbitrary computation. Corruption happens before any encoding. Thus the lower bounds cannot be removed by storing more of the corrupted input or by improving an internal representation. Conversely, the upper rules use counts directly; their performance does not rely on a restricted neural architecture.

The principal distinction is between protecting and not protecting the terminal pair. With protection, a logarithmic accumulation of rare informative suffixes determines the corruption price. Binary sources have a different cutoff from sources with three or more symbols. Without protection, even a known first-order source can leave irreducible uncertainty about the true final symbol after one erasure. The latter is uncertainty about the prediction target itself, not merely about estimating an unknown transition probability.

\paragraph{Relation to earlier work.}
Prediction from a Markov trajectory, and lower bounds based on a common prior and waiting-time structure, have established precedents; see Han, Jana and Wu \citep{han2022}. We do not claim that those general ideas originate here. Robust Markov inference also has a substantial existing literature. Ghosh \citep{ghosh2021} studies parametric inference for stationary ergodic finite chains, including contamination of the transition matrix by rowwise mixtures. That contamination acts on the transition law, rather than replacing a fixed number of entries in a realized word. Zuckerman \citep{zuckerman2026} studies robustness of Markov dynamics; its abstract concerns limiting frequencies and mixing, hitting and cover times for reversible chains. Bhatt et al.\ \citep{bhatt2022} study contaminated mean estimation under moment assumptions. These comparisons distinguish the declared experiments, not establish priority for the present rates. Our reading of Ghosh covers the abstract, introduction and Section 2.3; our comparisons with Zuckerman and Bhatt et al.\ are abstract-level only. The transport-coupling work of O'Connor, McGoff and Nobel \citep{oconnor2022} was screened at abstract level for long-run stationary transition couplings, not audited as a source of the finite bridge theorem proved here.

\paragraph{Proof organization and dependency boundary.}
We first state the experiment, rates and statistical-order budget consequence. The technical sections give the new perturbation, erasure, binary transport and endpoint arguments in full. They import a clean actual-terminal add-one bound from the predecessor dossier (R011), a stationary clipped-occupation inequality (IU001), and binary weighted-occupation interfaces (BIF001 and BAI003 version 2). The arbitrary-start residual decomposition is given explicitly, following AU001. These inherited interfaces are stated with their domains at the point of use. Their full original proofs and reviews remain in the accompanying predecessor dossier, whose exact paths and hashes are listed in the source manifest. This manuscript is therefore not a claim to newly rederive the entire foundational clean-prediction chain. The preserved companion manuscript on finite information cuts addresses different, internal representation restrictions.

\section{Experiment and principal results}\label{sec:model}
Let $\mathcal X$ be an alphabet of fixed size $K\ge2$. A homogeneous second-order source has transition vectors $\theta_{ab}\in\Delta_K$ and an initial pair law $\mu$. Put
\[
 N=M+2,\qquad Y_j=(X_{j+1},X_{j+2}),\quad 0\le j\le M,
 \qquad P((a,b),(b,c))=\theta_{ab}(c).
\]
All other pair-kernel entries are zero. Departures $j=0,\ldots,M-1$ give precisely $M$ observed triples. The target is $\theta_{Y_M}$, the conditional distribution of the unobserved next symbol.

For $0\le\beta\le1$, the class $\mathcal D_2(\beta)$ consists of tables admitting a stationary pair law $\Pi$ with
\begin{equation}\label{eq:certificate}
 \Pi P=\Pi,\qquad P^2(u,v)\ge\beta\Pi(v)\quad\text{for every pair }u,v\in\mathcal X^2.
\end{equation}
The certificate includes starting states outside stationary support. The qualifying $\Pi$ is existential and hidden; it need not be unique. In the stationary experiment the initial pair has such a qualifying law. In the arbitrary-initialization experiment $\mu$ is unrestricted, including unsupported starting pairs.

A protected corruption kernel $C$ is a stochastic map from $\mathcal X^N$ to itself with
\[
 d_H(x,z)\le k,\qquad (z_{N-1},z_N)=(x_{N-1},x_N)
 \quad\text{whenever } C(z\mid x)>0.
\]
It can inspect the complete realized word (an offline adversary) and use private randomness. It receives no unknown source table, stationary law, initial law, future symbol or hidden parameter as an input. Its algorithm may depend on public $K,M,k$. The predictor receives the entire $Z$, not an edit mask, attack seed or clean interior. It may use public $K,M,\beta,k$ and independent randomness. Loss is unconditional expected half-$\ell_1$ total variation,
\[
 \E\TV(A(Z),\theta_{Y_M}),\qquad
 \TV(p,q)=\tfrac12\sum_{c\in\mathcal X}|p(c)-q(c)|.
\]
Write
\begin{equation}\label{eq:risk}
 R_{\rm all}=\inf_A\sup_{(\theta,\mu)\in\mathcal D_2(\beta)}\sup_{C\in\mathcal C_k}
 \E\TV(A(C(X)),\theta_{Y_M}),
\end{equation}
and $R_{\rm stat}$ for the stationary restriction. The notation in the source supremum includes the unrestricted initial law only for $R_{\rm all}$. Nested suprema alone do not force the maximizing map to be identical at different sources. Our lower proofs establish the stronger fact that each witness family uses one common map across all of its unknown parameters, even if that map and family are known to the predictor. Separate lower floors may use different families and maps; their maximum, not an unsupported sum, is taken.

Throughout the main results $M\ge64$ and $0\le k\le M$ are integers. Define
\begin{align}
 s&=M^{-1/2},&
 e_K&=\begin{cases}\max\{\sqrt\beta,M^{-1}\},&K=2,\\
 \max\{\beta,M^{-1}\},&K\ge3,\end{cases}\label{eq:cutoff}\\
 F_K&=\begin{cases}0,&k=0,\\
 \dfrac{k}{M}\left[1+\log\left(1+\dfrac1{k e_K}\right)\right],&k\ge1,
 \end{cases}&
 H_K&=\min\{1,\max\{s,F_K\}\}.\label{eq:profile}
\end{align}
All logarithms are natural. A zero-budget expression is defined separately, never by evaluating a singular logarithm.

\begin{theorem}[Protected true terminal pair]\label{thm:main}
For every fixed $K\ge2$ there are positive constants $c_K,C_K$, independent of $M,k,\beta$, such that
\[
 c_K H_K\le R_{\rm stat}\le R_{\rm all}\le C_K H_K.
\]
One may take $c_K=2^{-15}$ for $K\ge3$ and $c_2=2^{-54}$. For $K\ge3$ the ordinary add-one rule applied to the corrupted terminal-row counts is both $\beta$-blind and $k$-blind and attains this rate. For $K=2$ a deterministic lower-envelope count rule attains it using public $k$, but not $\beta$. No unknown-$k$ binary adaptation claim is made.
\end{theorem}

The theorem is uniform in the stated finite domain, but its constants are not sharp in a growing alphabet. Stationary initialization does not eliminate the logarithmic price. The binary cutoff $\sqrt\beta$ comes from oriented terminal occupation and a legal binary witness; it is not a claim that all binary chains mix globally at that rate. For fixed positive $\beta$, both alphabets have risk of order $\max\{M^{-1/2},k/M\}$. Moving $\beta$ reveals the distinction.

\begin{theorem}[Editable endpoint, unchanged true target]\label{thm:endpoint-main}
Allow $C$ to edit any of the $N$ positions while retaining $d_H(X,Z)\le k$, and continue to target $\theta_{Y_M}$ at the original hidden pair. With $\delta=1-\beta$, set
\[
 H_{\rm end}=\min\{1,\max\{s,k/M,\delta\1_{\{k\ge1\}}\}\}.
\]
The two initialization risks in this experiment satisfy
\[
 H_{\rm end}/96\le R_{{\rm end},{\rm stat}}
 \le R_{{\rm end},{\rm all}}\le C_K^{\rm end}H_{\rm end},
 \quad C_K^{\rm end}=\max\{2^{13}K^{5/2},\sqrt K/2+3\}.
\]
For $k\ge1$, pooling corrupted symbols at positions $3,\ldots,N$ gives the upper bound; for $k=0$, use clean selected-row add-one.
\end{theorem}

In particular, a fixed $\beta<1$ and a single allowed endpoint edit preclude consistency for this target. A predictor asked instead for the law at the corrupted pair faces a different experiment. Revealing a mask or seed also changes the information available and is not covered here.

\section{How many replacements preserve statistical order?}\label{sec:tolerance}
The rate admits a useful explicit inverse. Define
\[
 f_K=1+\log\left(1+\frac1{e_K\sqrt M}\right),
 \qquad T_K=\frac{\sqrt M}{f_K}.
\]
This inverse is only an algebraic consequence of Theorem~\ref{thm:main}, not a further empirical result or a clean-reconstruction threshold.

\begin{corollary}[Statistical-order tolerance]\label{cor:tolerance}
On the domain of Theorem~\ref{thm:main}, $1\le T_K\le\sqrt M$.
Every integer $0\le k\le\lfloor T_K\rfloor$ satisfies $F_K\le2s$.
Conversely, if $F_K\le A s$ for a fixed $A>0$ and $k\ge1$, then, with $A'=\max\{A,1\}$,
\[
 k\le A A' T_K.
\]
Consequently the largest budget retaining risk of statistical order is of order $T_K$, with separate constants in sufficiency and necessity. For the endpoint-unprotected experiment a positive budget retains statistical order exactly in the order sense $k=O(\sqrt M)$ and $1-\beta=O(M^{-1/2})$.
\end{corollary}
\begin{proof}
Write $e=e_K$, $x=1/(e\sqrt M)$, $f=1+\log(1+x)$ and $T=\sqrt M/f$. Since $e\ge1/M$, $0<x\le\sqrt M$. For $r\ge8$, $1+\log(1+r)\le r$ (check it at $8$ and differentiate). Thus $1\le f\le\sqrt M$, proving the range of $T$.
For real $u>0$, the function $u[1+\log(1+1/(eu))]$ is increasing: its derivative is $1+\log(1+v)-v/(1+v)>0$, where $v=1/(eu)$. At $u=T$,
\[
 F(T)/s=\frac{1+\log(1+fx)}{f}
 \le\frac{f+\log f}{f}\le2.
\]
Here $1+fx\le f(1+x)$ and $\log f\le f$. Monotonicity gives the asserted integer sufficiency, including $k=0$ by definition. Since $T\ge1$, $\lfloor T\rfloor\ge T/2$, so rounding does not remove the order-sized positive budget.

For necessity, $F_K\ge k/M$ first gives $k\le A\sqrt M\le A'\sqrt M$. Therefore
\[
 1+\log(1+1/(ek))\ge1+\log(1+x/A')\ge f/A'.
\]
The last step uses concavity of $\log(1+x)$ at the factor $1/A'\in(0,1]$ and $1\ge1/A'$. Combining with $F_K\le As$ gives $k\le AA'T$. These are not a same-multiplier ``if and only if.'' More generally, for $0<c\le1$, monotonicity and $1+fx/c\le(f/c)(1+x)$ give $F(cT)\le c[2+\log(1/c)]s$, permitting any prescribed positive fixed profile multiplier by decreasing $c$ and rounding down; when the resulting threshold is below one, only zero budget is guaranteed.

For a sequence with $M\to\infty$, a risk bounded by a fixed multiple of $s$ eventually lies below the cap in Theorem~\ref{thm:main}. Its lower bound therefore forces $F_K=O(s)$; the upper gives the converse. Finitely many remaining $M$ are absorbed into constants. Finally Theorem~\ref{thm:endpoint-main} directly gives the two stated conditions for a positive unprotected budget. Zero budget is a separate clean case and imposes no condition on $1-\beta$.
\end{proof}

At $\beta=1/M$, the binary value has $e_2=M^{-1/2}$ and $f_2=1+\log2$, so $T_2$ is of order $\sqrt M$. For $K\ge3$, $e_K=M^{-1}$ and $f_K=1+\log(1+\sqrt M)$, giving order $\sqrt M/\log M$. At $\beta=0$ both cutoffs equal $M^{-1}$ and both tolerances have the latter order. In the unprotected experiment, $\beta=1/M$ instead leaves a constant-order one-edit obstruction as $M$ grows.

\section{Scope of the synthesis}
The protected upper and lower directions have different mechanisms. For $K\ge3$, direct histogram stability is combined with terminal occupation, while exact exchangeability erases the active suffix arrangement. For binary sources, a pessimistic designated-exit estimator exploits a weighted occupation inequality. Its lower witness has variable-length cycles, so a cheap modification with the wrong output marginal is not enough. Entropy transport supplies a weaker bridge approximation, and a finite scalar argument shows that its genuine horizon cutoff does not change the final order.

The binary lower is an existence result for a common stochastic map assembled from finite rational transport tables. It makes no efficient sampling claim. The upper count rule has a straightforward finite rational description, but no assertion about a Transformer's bounded depth, width, attention budget, training dynamics or learnability follows. These resources require separate theorems with the same input and fault permissions. Likewise, offline history-adaptive kernels, causal corruption and visible-mask models are not interchangeable here.

The technical proofs below imply Theorems~\ref{thm:main} and~\ref{thm:endpoint-main} by taking maxima of the separately established lower floors and using the stated upper rules. All genuinely new ingredients are included; the explicitly named inherited interfaces retain their original evidence status. Independent LLM reviews and finite diagnostics in the dossier are corroborating mathematical checks, not formal verification or external peer review.

\appendix
\section{Protected endpoints with at least three symbols}
\label{sec:nonbinary}
The logarithm in the corruption rate comes from terminal states with few
previous visits.  It is not an artifact of a compressed representation.
We first bound the perturbation of a direct full-history estimator, and
then construct an erasure channel whose \emph{entire output word} has only
a controlled amount of parameter information.
Throughout this section $K\geq3$, $M\geq64$, $0\leq\beta\leq1$, and
$k\in\{0,\ldots,M\}$.  Put $b_0=\max\{\beta,M^{-1}\}$,
$s=M^{-1/2}$, and
\[
 F_3=\begin{cases}
 0,&k=0,\\
 \displaystyle\frac{k}{M}\left[1+\log\left(1+\frac1{b_0k}\right)\right],&k\geq1.
 \end{cases}
\]

\subsection{Pathwise stability of corrupted counts}
Write $C_{abc}$ for the number of clean triples $(a,b,c)$ ending at
positions $3,\ldots,M+2$, and $C'_{abc}$ for the corresponding corrupted
counts.  A replacement affects at most three triple locations.  At each
location, the old and new indicator vectors differ in $\ell_1$ norm by at
most two.  Consequently, even with adjacent or simultaneous edits,
\begin{equation}\label{eq:nb-histogram}
 \|C'-C\|_1\leq6k.
\end{equation}
The terminal context $S=(X_{M+1},X_{M+2})$ is preserved.  Set
$J=\sum_cC_{Sc}$ and let
\[
 q_c=\frac{C_{Sc}+1}{J+K},\qquad
 q'_c=\frac{C'_{Sc}+1}{\sum_dC'_{Sd}+K}.
\]
For nonnegative vectors $u,v$ of positive masses $A,B$, insertion of
$v/A$ and the inequality $|A-B|\leq\|u-v\|_1$ give
\[
 \operatorname{TV}(u/A,v/B)
 \leq\tfrac12\bigl(\|u-v\|_1/A+|A-B|/A\bigr)
 \leq\|u-v\|_1/A.
\]
Apply this to the two smoothed rows, using the \emph{clean} mass in the
denominator, to obtain the deterministic bound
\begin{equation}\label{eq:nb-row-stability}
 \operatorname{TV}(q',q)\leq\min\{1,6k/(J+K)\}.
\end{equation}
This argument allows arbitrary changes to row membership.  It does not
condition row successors on the selected endpoint or treat them as IID.

We use the following inherited analytic interfaces.  They are stated here
to delimit the dependence on the predecessor dossier; the new corruption
argument does not re-prove its clean statistical theorem or its stationary
reverse-return analysis.  The R011 ORIGINAL rule is exactly the clean
add-one vector $q$ above.  Its uniform consequence applies for every
$K\geq2$, integer $M\geq1$, $\beta\in[0,1]$, every qualifying source,
and every initial pair law (including stationary initialization):
\begin{equation}\label{eq:nb-clean-interface}
 \mathbb E\operatorname{TV}(q,\theta_S)
 \leq2^{13}K^{5/2}M^{-1/2}.
\end{equation}
R011 gives a finer parameter-dependent bound, but its full collection
of scales is not needed here: the displayed uniform consequence is the
only clean upper interface used in this manuscript.  Its $K=2$ domain
also covers the endpoint result in Section~\ref{sec:endpoint}.

For a pair-chain path $Y_0,\ldots,Y_m$, define its actual-terminal
occupation count by
$J_m=\sum_{i=0}^{m-1}\mathbf1\{Y_i=Y_m\}$.
The IU001 stationary interface, for $\beta>0$, $m\geq1$, $A>0$, and
integer $\kappa\geq1$, is
\begin{equation}\label{eq:nb-IU-interface}
 \mathbb E_{\Pi}\min\{1,A/(J_m+\kappa)\}
 \leq\frac{dA}{m}\left[17+\log\left(1+\frac{480}{\beta A}\right)\right],
 \qquad d=|\operatorname{supp}\Pi|\leq K^2.
\end{equation}
This is a terminal-query bound under the two-step all-state certificate,
not an independent stationary-query identity.  The AU001 universal
interface, valid for every initial law, every $\beta\in[0,1]$, and the
same domains of $m,A,\kappa$, is
\begin{equation}\label{eq:nb-AU-interface}
 \mathbb E_{\nu}\min\{1,A/(J_m+\kappa)\}
 \leq\Psi\left(\frac{K^4A}{m+1}\right),\qquad
 \Psi(x)=\begin{cases}x(1+\log(1/x)),&0<x<1,\\0,&x=0,\\1,&x\geq1.
 \end{cases}
\end{equation}
For clarity, the short universal occupation argument can be reproduced.
For a chain on $d$ states and integer $w\geq1$, put
$b_m(y)=\mathbb P_y(J_m<w)$ and $B_m=\max_y b_m(y)$.
The count over a whole path dominates the count over an ending suffix,
so the Markov property gives $B_{m+1}\leq B_m$.  On each path,
$\sum_{t=0}^m\mathbf1\{J_t<w\}\leq dw$: each state can be encountered
with fewer than $w$ previous visits at most $w$ times.  Thus
\[
 (m+1)B_m\leq\sum_{t=0}^m B_t
 \leq\sum_y\sum_{t=0}^m b_t(y)\leq d^2w.
\]
Layer cake proves \eqref{eq:nb-AU-interface}.  Indeed, for $0<t<1$ the
strict threshold $J_m<A/t-\kappa$, when nonempty, is represented by
$w=\lceil A/t-\kappa\rceil\leq A/t$; integrate
$\min\{1,K^4A/((m+1)t)\}$.  This also explains why the smoothing
integer and the strict inequality matter.

\subsection{An exact suffix mixture for arbitrary initialization}
Assume first $\beta>0$ and $k\geq1$.  Let
\[
 u=\lfloor M/2\rfloor,\quad m=\lceil M/2\rceil,\quad
 r=\lfloor M/4\rfloor,\quad \eta=(1-\beta)^r.
\]
For $0<\beta<1$, the matrix
$R=(P^2-\beta\mathbf1\Pi)/(1-\beta)$ is stochastic on \emph{all}
pair states, and $\Pi R=\Pi$.  Expanding the powers of
$P^2=\beta\mathbf1\Pi+(1-\beta)R$ gives
\begin{equation}\label{eq:nb-suffix-mixture}
 \mu P^u=(1-\eta)\Pi+\eta\nu,
 \qquad \nu=\mu R^rP^{u-2r}.
\end{equation}
Extending each initial distribution with the ordinary transition matrix
$P$ gives the same exact mixture for the \emph{entire suffix path}
$Y_u,\ldots,Y_M$.  In particular, the residual suffix uses $P$, not $R$.
For $\beta=1$, $P^2=\mathbf1\Pi$ and the residual term is absent.
The full terminal count dominates the suffix terminal count pathwise.
Apply \eqref{eq:nb-IU-interface} and \eqref{eq:nb-AU-interface} with
$A=6k$, $\kappa=K$, and $m\geq M/2$:
\begin{align}\label{eq:nb-suffix-bound}
 \mathbb E\min\{1,6k/(J+K)\}
 &\leq12K^2\frac{k}{M}
 \left[17+\log\left(1+\frac{80}{\beta k}\right)\right]
 +\eta\Psi(12K^4k/M).
\end{align}
No synthetic regeneration event is supplied to the estimator; this
mixture is only an analysis of its existing path distribution.

For $\beta\geq1/M$, write $z=M/k\geq1$, $t=M\beta\geq1$,
$D=12K^4$, and $G=1+\log(1+z/t)$, so $F_3=G/z$.
We have $\eta\leq1$ and $\eta\leq e\exp(-t/4)$.
If $z<D$, then $\eta\Psi(D/z)\leq1\leq DG/z$.
If $z\geq D$, use
$\log z\leq\log(1+z/t)+\log t$ to obtain
\[
 \eta\Psi(D/z)\leq\frac Dz\eta(1+\log z)
 \leq\frac Dz\{G+e\exp(-t/4)\log t\}
 \leq\frac{5DG}{z}.
\]
The last elementary bound uses $e\exp(-t/4)\log t\leq4$ for $t\geq1$.
Also $\log(1+80x)\leq\log80+\log(1+x)$.
For $0\leq\beta<1/M$, apply \eqref{eq:nb-AU-interface} directly to
the full path.  Splitting at $z=6K^4$ proves
\[
 \Psi\left(\frac{6K^4k}{M+1}\right)
 \leq\Psi(6K^4/z)
 \leq\frac{6K^4}{z}[1+\log(1+z)]=6K^4F_3.
\]
Combining these estimates with \eqref{eq:nb-row-stability} and
\eqref{eq:nb-clean-interface} gives
\begin{equation}\label{eq:nb-final-upper}
 \mathbb E\operatorname{TV}(q',\theta_S)
 \leq\min\left\{1,2^{13}K^{5/2}s+
 [12K^2(17+\log80)+60K^4]F_3\right\}.
\end{equation}
At $k=0$ the perturbation vanishes and only the clean interface is used.
The rule $q'$ itself uses neither $k$ nor $\beta$.

\subsection{Exact erasure of arrangements}
\begin{lemma}[Exchangeable count erasure]\label{lem:nb-erasure}
Let $V_1,\ldots,V_\ell$ be IID Bernoulli$(p)$ bits, set
$n=\sum_iV_i$, fix an integer $b\geq1$, and put
$r=b\lfloor n/b\rfloor$.  Choose uniformly $n-r$ of the one positions
and replace them by zero.  The output $Z$, conditional on $r$, is uniform
among weight-$r$ words, independently of $p$.  At most $b-1$ bits are
changed.  If two further IID Bernoulli$(p)$ bits $U$ are untouched, the
joint experiment $(Z,U)$ factors through $(r,U)$ followed by a
$p$-independent channel, with at most
$4(\lfloor\ell/b\rfloor+1)$ messages.
\end{lemma}
\begin{proof}
Conditional on $n$, the original word is uniform among its
$\binom\ell n$ arrangements.  A fixed output of weight $r$ is contained
in $\binom{\ell-r}{n-r}$ original one sets, and is chosen from each one
set with probability $1/\binom nr$.  Hence
\[
 \mathbb P(Z=z\mid n)
 =\frac{\binom{\ell-r}{n-r}}{\binom\ell n\binom nr}
 =\frac1{\binom\ell r}.
\]
Averaging within a cell of the rounded count leaves this identity
unchanged.  Conditional on $p$, the two untouched bits are independent
of the first $\ell$ bits and the erasure randomization.  Thus the same
uniform-arrangement channel works conditional on $(r,U)$.
The identity includes $\ell=0$ and the incomplete last count cell.
\end{proof}

\begin{lemma}[Finite-message Bayes loss]\label{nb:finite-message}
If a real parameter is uniform on an interval of length $h>0$ and an
arbitrary stochastic channel emits at most $D$ messages, every
independently randomized real-valued decoder has expected absolute loss
at least $h/(4D)$.
\end{lemma}
\begin{proof}
Fix its independent decoding seed, obtaining at most $D$ possible output
points.  Their radius-$t$ neighborhoods have prior mass at most $2Dt/h$.
Whatever stochastic encoding is used, the event of error at most $t$
requires that the parameter lie in that union.  Integrate
$1-2Dt/h$ over $0\leq t\leq h/(2D)$ and then average the seed.
\end{proof}
In Lemma~\ref{lem:nb-erasure}, a full-word decoder can be composed with
the parameter-independent uniform-arrangement channel, so it is a
randomized decoder of the finite message.  This is not an assumption
that the observation is only a histogram.  The protected pair is part
of the message, and the prior is not asserted to remain uniform after
conditioning on that pair.  Neither the actual edit mask nor the actual
attack seed is disclosed.

\subsection{Two IID floors common to all mixing classes}
\begin{lemma}\label{nb:iid-floors}
For every $K\geq2$, $M\geq64$, and $\beta\in[0,1]$, the stationary
protected-endpoint minimax risk is at least $s/64$.  For $k\geq1$ it is
also at least $k/(96M)$.  Each lower bound is realized by a stationary
IID binary family and a common corruption kernel across that family.
\end{lemma}
\begin{proof}
Choose two named symbols $C_0,C_1$ and set every source row, including
unused contexts, to $(1-p)\delta_{C_0}+p\delta_{C_1}$.  Its two-step
pair matrix has every row equal to the stationary pair law, so it
satisfies every certificate under consideration.
For the statistical floor use identity corruption and
$p_\pm=1/2\pm1/(16\sqrt N)$, $N=M+2$.
The one-sample divergence is at most $32/(256N)$, hence the product
KL divergence is at most $1/8$ and Pinsker's inequality bounds product
TV by $1/4$.  The target laws are separated in TV by $1/(8\sqrt N)$.
Nearest-target testing and the triangle inequality give average loss at
least one quarter of this separation times $1-1/4$, namely
$3/(128\sqrt N)\geq s/64$.

For the corruption floor use one prior uniform on $[1/4,3/4]$ and
$b=k+1$.  On all-binary input words apply Lemma~\ref{lem:nb-erasure}
to the first $M$ bits and leave the last two untouched; on every other
input use identity.  This defines a common map on the entire alphabet
with at most $k$ edits.  The full output has at most
$D=4(\lfloor M/b\rfloor+1)$ informative messages.
Since TV loss dominates the absolute error in the $C_1$ coordinate,
Lemma~\ref{nb:finite-message} gives
\[
 R_{\mathrm{stat}}^{\mathrm{prot}}
 \geq\frac1{32(\lfloor M/b\rfloor+1)}
 \geq\frac b{32(M+b)}\geq\frac{k}{96M},
\]
where $b\leq M+1\leq2M$.  The statistical and corruption witnesses
need not coincide: their lower bounds are combined by a maximum.
\end{proof}

\subsection{Stationary rare suffixes and the full logarithm}
Suppose $b_0\leq1/16$.  Put $a=1/(2M)$ and $e=b_0-a$.
For $p\in[1/4,3/4]$, $p_0=1-p$, $p_1=p$, use symbols $W,C_0,C_1$
and define a first-order chain by
\begin{align*}
 T(W,W)&=1-a,& T(W,C_i)&=ap_i,\\
 T(C_i,W)&=e,&T(C_i,C_j)&=(1-e)p_j.
\end{align*}
Its stationary symbol law is
\[
 \pi(W)=e/b_0,\qquad\pi(C_i)=ap_i/b_0.
\]
Give unused symbols zero stationary mass and set every unused starting
row to $\pi$.  Direct substitution verifies $\pi T=\pi$.
The coordinate inequalities $T(y,v)\geq b_0\pi(v)$ hold for every
starting symbol $y$, including unused ones: the $W$-to-active and
active-to-$W$ ratios equal $b_0$, while the other inequalities reduce
to $b_0\leq1$.  Set $\theta_{xy}(v)=T(y,v)$ on \emph{every} pair and
$\Pi(u,v)=\pi(u)T(u,v)$.  Then $\Pi$ is stationary and
\begin{equation}\label{eq:nb-allstate-witness}
 P^2((x,y),(u,v))=T(y,u)T(u,v)
 \geq b_0\Pi(u,v)\geq\beta\Pi(u,v)
\end{equation}
on all pair states, not merely the stationary support.
Initialize at this stationary pair law, which is not $\pi\otimes\pi$.

Let $L=\lfloor1/(4b_0)\rfloor$, so $4\leq L\leq M/4$.
For $2\leq j\leq L$, let $E_j$ be the event that the complete word is
$W^{M+2-j}$ followed by $j$ active letters.  These disjoint events have
probabilities
\[
 w_j=\Pi(W,W)(1-a)^{M-j}a(1-e)^{j-1}.
\]
Here $\Pi(W,W)=(1-a/b_0)(1-a)\geq1/4$ is independent of $p$.
Bernoulli's inequality gives
$(1-a)^{M-j}\geq1/2$ and
$(1-e)^{j-1}\geq1-b_0L\geq3/4$.  Thus
\begin{equation}\label{eq:nb-event-mass}
 w_j\geq3/(64M)\geq1/(32M).
\end{equation}
Conditional on $E_j$, the $j$ active bits are IID Bernoulli$(p)$.
The same uniform prior on $p$ is retained after this event, before the
endpoint labels are observed.

Define a single kernel on all words as follows.  If its input is
$W^{N-j}$ followed by $j$ active symbols with $2\leq j\leq M$, round
the number of $C_1$ symbols among the first $j-2$ active symbols to a
multiple of $b=k+1$, using the uniform deletion rule.  Preserve all
markers and the final pair.  Use identity on every other word.
This kernel depends on neither $p$ nor $\beta$, preserves the visible
suffix length $j$, and makes at most $k$ replacements in allowed
positions on every input.  It is the same kernel for every suffix and
every parameter value.

At all active terminal pairs the actual future law is
\begin{equation}\label{eq:nb-true-future}
 e\delta_W+(1-e)[(1-p)\delta_{C_0}+p\delta_{C_1}].
\end{equation}
The return mass $e$ must not be omitted merely because no return
occurred in the selected past.  For every prediction $q$,
$\operatorname{TV}(q,\theta_S)\geq|q(C_1)-(1-e)p|$.
Apply Lemmas~\ref{lem:nb-erasure} and \ref{nb:finite-message} to
$q(C_1)/(1-e)$, with interval length $1/2$.  Since $1-e\geq1/2$,
the conditional Bayes loss is at least
\[
 \frac1{64[1+(j-2)/b]}.
\]
The decoder may be granted the value of $j$, which is already visible.
The output arrangement adds no information beyond the finite message.
Sum the losses over the disjoint events, without hiding observations
outside them:
\begin{align}\label{eq:nb-harmonic-lower}
 R_{\mathrm{stat}}^{\mathrm{prot}}
 &\geq\frac1{2048M}\sum_{\ell=0}^{L-2}\frac1{1+\ell/b}
 \geq\frac b{2048M}\log\left(1+\frac{L-1}{b}\right).
\end{align}
The last inequality integrates a decreasing function over $[0,L-1]$.
Since $L-1\geq1/(8b_0)$, the function
$b\log(1+c/b)$ is increasing in $b>0$, and
$\log(1+x/8)\geq\log(1+x)/8$, it follows that
\begin{equation}\label{eq:nb-log-lower}
 R_{\mathrm{stat}}^{\mathrm{prot}}
 \geq\frac{k}{16384M}\log\left(1+\frac1{b_0k}\right).
\end{equation}
The same prior and same kernel were used across all $j$; choosing a
different source separately at each suffix would not justify this sum.

For $b_0\leq1/16$, maximize \eqref{eq:nb-log-lower} and the IID
$k/(96M)$ floor to obtain $F_3/32768$.  For $b_0>1/16$,
$\log(1+1/(b_0k))\leq\log17<3$, so $F_3\leq4k/M$ and the IID floor
alone is stronger than $F_3/32768$.  Maximizing with the statistical
floor gives
\begin{equation}\label{eq:nb-final-lower}
 R_{\mathrm{stat}}^{\mathrm{prot}}
 \geq2^{-15}\min\{1,\max(s,F_3)\}.
\end{equation}
For $k=0$, use only the statistical floor.  Together with
\eqref{eq:nb-final-upper} and inclusion of stationary initializations
among arbitrary ones, this proves the claimed nonbinary profile.

\subsection{The zero-certificate mechanism without a mixing cutoff}
The preceding construction already covers $\beta=0$, but a simpler
version illuminates the logarithm and records the original stationary
strengthening.  This paragraph is valid for every $M\geq2$, not only
$M\geq64$.
For arbitrary initialization, start at $WW$, set $\alpha=1/M$, and use
\[
 T(W,W)=1-\alpha,\quad T(W,C_i)=\alpha p_i,
 \quad T(C_i,C_j)=p_j.
\]
Unused rows can equal $\delta_W$.  A qualifying stationary law is
supported on the active IID class, so the source belongs to
$\mathcal D_2(0)$ despite its nonstationary start.  For
$2\leq j\leq M$, the event $E_j$ above has mass
$\alpha(1-\alpha)^{M-j}\geq1/(4M)$, independent of $p$.
The future active law is exactly Bernoulli$(p)$.
The same marker erasure gives conditional risk at least
$1/[32(1+(j-2)/(k+1))]$.  Hence, for $b=k+1$,
\[
 R_{\mathrm{all}}^{\mathrm{prot}}
 \geq\frac b{128M}\log\left(1+\frac{M-1}{b}\right)
 \geq\frac1{512}\Psi(k/M),\qquad k\geq1.
\]
For the last inequality, use $M-1\geq M/2$, monotonicity in $b$, and
$\log(1+t/2)\geq(1+\log t)/4$ for $t=M/k\geq1$.
The latter follows by checking $t=1$ and differentiating.

For a stationary version set $a=e=1/M$ in the recurrent marker chain.
Now $\pi(W)=1/2$, $\pi(C_i)=p_i/2$, and the initialization is
$\Pi(u,v)=\pi(u)T(u,v)$.  For every $2\leq j\leq M$,
\[
 \mathbb P_p(E_j)
 =\frac{(1-1/M)^M}{2M}\geq\frac1{8M}.
\]
The equality includes the initial $WW$ factor; it is independent of
both $p$ and $j$.  The lower bound follows from
$(1-1/M)^M\geq1/4$ for $M\geq2$.
The true future law retains its return mass $e$, reducing the
conditional lower by at most a factor two.  The same summation yields
$\Psi(k/M)/2048$.  Thus transience is not needed for the logarithmic
obstruction.  No concentration cutoff such as $j=O(k^2)$ is needed:
exact arrangement erasure works for every suffix length.

\section{When the original terminal context is not protected}
\label{sec:endpoint}
We now change exactly one part of the experiment.  The adversary may
replace any of the $N=M+2$ positions, including the final two, subject
to the same strict budget $k\leq M$.  The prediction target remains
$\theta_{X_{N-1},X_N}$ at the \emph{original hidden} terminal pair.
It is not the row indexed by the corrupted endpoint.  Write
$R_{\mathrm{stat}}^{\mathrm{end}}$ and
$R_{\mathrm{all}}^{\mathrm{end}}$ for these risks and put
\[
 \delta=1-\beta,\qquad
 H_{\mathrm{end}}=\min\{1,\max(s,k/M,\delta\mathbf1\{k\geq1\})\}.
\]

\begin{proposition}\label{prop:endpoint}
For $K\geq2$, $M\geq64$, $0\leq\beta\leq1$, and integer
$0\leq k\leq M$,
\[
 H_{\mathrm{end}}/96\leq R_{\mathrm{stat}}^{\mathrm{end}}
 \leq R_{\mathrm{all}}^{\mathrm{end}}
 \leq C_K^{\mathrm{end}}H_{\mathrm{end}},
 \qquad
 C_K^{\mathrm{end}}=\max\{2^{13}K^{5/2},\sqrt K/2+3\}.
\]
For $k=0$ the clean selected-row add-one rule suffices.  For $k\geq1$
the pooled empirical distribution of corrupted symbols at positions
$3,\ldots,N$ suffices.  Neither branch requires $\beta$.
\end{proposition}

\subsection{What the all-state certificate says about each row}
Let $\pi(c)=\sum_d\Pi(c,d)$.  Summing
$P^2(u,(c,d))\geq\beta\Pi(c,d)$ over $d$ gives
\begin{equation}\label{eq:end-row-domination}
 \theta_u(c)\geq\beta\pi(c)
 \quad\hbox{for every pair $u$ and symbol $c$}.
\end{equation}
Indeed the first coordinate after two pair transitions is precisely the
next generated symbol from $u$.  The conclusion includes contexts
outside stationary support.  For $\beta<1$, write
$\theta_u=\beta\pi+\delta r_u$, where $r_u$ is a probability vector.
For $\beta=1$, coordinatewise domination and unit mass imply
$\theta_u=\pi$.  In every case,
\begin{equation}\label{eq:end-row-tv}
 \operatorname{TV}(\theta_u,\pi)\leq\delta.
\end{equation}
At $\beta=1$ all generated symbols from position 3 onward are IID
$\pi$, even for an arbitrary initial pair; no assertion about that
initial pair being IID is needed.

\subsection{A pooled martingale upper bound}
For $k\geq1$, let
$\bar p=M^{-1}\sum_{t=3}^N e_{X_t}$ be the clean generated-symbol
empirical distribution.  With the natural filtration, define
$m_t=\theta_{X_{t-2},X_{t-1}}$ and $D_t=e_{X_t}-m_t$.
These are vector martingale differences.  Orthogonality at distinct
times gives
\[
 \mathbb E\left\|\frac1M\sum_{t=3}^ND_t\right\|_2^2
 =\frac1{M^2}\sum_{t=3}^N\mathbb E[1-\|m_t\|_2^2]
 \leq\frac1M.
\]
Cauchy--Schwarz and Jensen imply that the expected half-$\ell_1$ norm
of this martingale average is at most $\sqrt K/(2\sqrt M)$.
By \eqref{eq:end-row-tv} and convexity, the average $M^{-1}\sum_tm_t$
is within TV $\delta$ of $\pi$ on every path, as is the true terminal
next law.  A pathwise triangle inequality therefore yields
\[
 \mathbb E\operatorname{TV}(\bar p,\theta_{X_{N-1},X_N})
 \leq\frac{\sqrt K}{2\sqrt M}+2\delta.
\]
Let $\bar p'$ pool the corrupted symbols at positions $3,\ldots,N$.
Each changed symbol alters this empirical distribution by at most
$1/M$ in TV, so $\operatorname{TV}(\bar p',\bar p)\leq k/M$ pathwise.
Edits to the first two positions contribute nothing to this bound.
Thus
\begin{equation}\label{eq:end-upper}
 \mathbb E\operatorname{TV}(\bar p',\theta_{X_{N-1},X_N})
 \leq\min\{1,\sqrt K s/2+2\delta+k/M\}.
\end{equation}
This proof does not reconstruct the endpoint, appeal to mixing-time
concentration, or substitute an independent query for the actual one.
For $k=0$, use \eqref{eq:nb-clean-interface} instead: the pooled bias
$2\delta$ is only an upper-bound artifact in the clean problem.
These two branches prove the upper in Proposition~\ref{prop:endpoint}.

\subsection{A one-symbol erasure lower bound even for a known source}
Assume $k\geq1$.  Set $\alpha=\min\{\delta/2,1/4\}$ and consider the
symmetric first-order binary chain
\[
 T(0,0)=T(1,1)=\tfrac12+\alpha,
 \qquad T(0,1)=T(1,0)=\tfrac12-\alpha.
\]
Its stationary symbol law is uniform on $\{0,1\}$.  For $K>2$, set
unused starting rows to that same binary uniform law and all unused
stationary masses to zero.  Extend to every second-order row by
$\theta_{xy}(c)=T(y,c)$ and initialize at
$\Pi(u,v)=\pi(u)T(u,v)$.  Every starting symbol satisfies
$T(y,u)\geq(1-2\alpha)\pi(u)$, including unused symbols, so
\[
 P^2((x,y),(u,v))=T(y,u)T(u,v)
 \geq(1-2\alpha)\Pi(u,v)\geq\beta\Pi(u,v).
\]
The witness therefore obeys the required all-state certificate.

Use the deterministic common attack that sets only the last symbol to
$0$ and leaves the rest unchanged.  It makes at most one replacement
on every word and requires no source parameters.  Grant the estimator
the entire source law and its stationary initialization.  The output
then reveals the true prefix through $X_{N-1}$ but not $X_N$.
Conditionally on that prefix, the probabilities of the two possible
last bits are $1/2+\alpha$ and $1/2-\alpha$ in some order.  The two
possible true next-law probabilities of symbol $1$ are
$1/2-\alpha$ and $1/2+\alpha$.
For any real prediction, the least posterior absolute error is the
smaller last-bit probability times the gap between these target values:
\[
 (\tfrac12-\alpha)\,2\alpha=\alpha-2\alpha^2.
\]
To check this weighted-median formula, the objective is linear below,
between, and above the two target values, and its minimum occurs at
the more likely endpoint.  Randomization cannot lower a linear average
of deterministic risks.  TV loss dominates the absolute error in the
symbol-$1$ coordinate, including predictions assigning mass to unused
symbols.  The bound holds after every prefix and hence unconditionally.

If $\delta\leq1/2$, the loss is
$\delta(1-\delta)/2\geq\delta/4$; if $\delta\geq1/2$, it is
$1/8\geq\delta/8$.  Therefore
\begin{equation}\label{eq:end-one-edit}
 R_{\mathrm{stat}}^{\mathrm{end}}\geq\delta/8,
 \qquad k\geq1.
\end{equation}
This obstruction is uncertainty about a latent final symbol in a
\emph{known} source, rather than uncertainty in estimating its rows.
The model does not reveal a replacement indicator: for this particular
attack that indicator would itself reveal the erased binary symbol.
No conclusion about all possible visible-mask models follows.

\subsection{Assembly and scope of the endpoint contrast}
Both IID floors in Lemma~\ref{nb:iid-floors} remain valid.  Their
identity and first-$M$-symbol erasure kernels are allowed in the enlarged
corruption class, even though they do not use endpoint edits.
For $k\geq1$, maximize $s/64$, $k/(96M)$, and
\eqref{eq:end-one-edit}; this gives
$\max(s,k/M,\delta)/96$, and hence the required capped lower.
For $k=0$, only the statistical floor is needed.
All lower witnesses are stationary, and the inclusion of stationary
initializations among arbitrary ones completes the proof.

In particular, at any fixed $\beta<1$, one allowed endpoint edit
prevents consistency for the original-terminal target.  At $\beta=1$
all source rows coincide and the endpoint ambiguity disappears.
For positive budgets, statistical-order risk is possible only when
$1-\beta=O(M^{-1/2})$ and $k=O(\sqrt M)$, and these two conditions
are also sufficient up to fixed-$K$ constants by the proved bounds.
This comparison concerns full-history statistical prediction.  It does
not assert an online adversarial equivalence, a visible-mask result,
a neural-depth requirement, or a training guarantee.

\section{The binary upper bound}
\label{bin:upper}
The distinction between two symbols and three or more symbols enters through a weighted occupation functional, not through a uniform improvement of the chain's mixing time. This section gives a fixed, publicly oriented estimator whose corruption error is bounded by that functional. Write
\[
 Y_j=(X_{j+1},X_{j+2}),\quad 0\le j\le M,
 \qquad J_u=\sum_{j=0}^{M-1}\mathbf 1\{Y_j=u\}.
\]
All row counts below use these $M$ departures, and never count the terminal query as a departure. Throughout this section $K=2$, $M\ge64$, and
\[
 s=M^{-1/2},\qquad e_*=\max\{\sqrt\beta,M^{-1}\}.
\]
For the corruption profile set
\[
 F_2=\frac{k}{M}\left[1+\log\left(1+\frac1{k e_*}\right)\right]
 \quad(k\ge1),\qquad F_2=0\quad(k=0).
\]

\subsection{A fixed orientation and a pessimistic exit estimate}
For a pair $u$, designate the successor $d(u)$ to be the complement of its first bit. In the order $00,01,10,11$, the designated successors are $1,1,0,0$. Their true probabilities, used only in the analysis, are
\[
 q=(a,b,e,t)=\bigl(\theta_{00}(1),\theta_{01}(1),
                    \theta_{10}(0),\theta_{11}(0)\bigr).
\]
The symbol $e$ here is a row probability, distinct from the cutoff $e_*$. The orientation does not require identifying a minority outcome or knowing the source.

At the preserved terminal pair $u=Y_M$, let $C'$ be the corrupted designated-successor count and $J'$ the corrupted row total. For $k\ge1$, set $B=6k$ and return a binary vector with designated coordinate
\begin{equation}
 \widehat q=\frac{(C'-B)_+}{J'+B+2}.
 \label{bin:estimator}
\end{equation}
Its other coordinate is $1-\widehat q$. At $k=0$, use ordinary add-one, with coordinates $(C'_{uv}+1)/(J'+2)$. These are different, explicitly specified branches. Both are deterministic and $\beta$-blind; the first uses the public budget $k$.

\begin{lemma}[Pathwise lower envelope]
\label{bin:envelope}
For $k\ge1$, write $C,J$ for the clean designated count and row total at $Y_M$, and put $A=24k+2$. Every admissible corrupted history satisfies
\begin{equation}
 \left|\widehat q-q_{Y_M}\right|
 \le 2\mathbf 1\{J>0\}\left|\frac CJ-q_{Y_M}\right|
       +\min\left\{q_{Y_M},\frac{A}{J+2}\right\}.
 \label{bin:pathloss}
\end{equation}
\end{lemma}
\begin{proof}
One changed symbol belongs to at most three counted triples. A replacement of one triple label changes the complete histogram in $\ell^1$ by at most two. Even with overlapping or adjacent edits, the histogram change is at most $6k=B$. Restriction to the selected binary row gives
\[
 |C'-C|+|(J'-C')-(J-C)|\le B.
\]
Endpoint protection is essential: the selected clean and corrupted row is the same. Set $n'=(C'-B)_+$, $D=J+2$, and $D'=J'+B+2$. Then
\[
 0\le n'\le C,\quad D'\ge D,\quad C-n'\le2B,
 \quad 0\le D'-D\le2B.
\]
Consequently
\[
 0\le \frac CD-\frac{n'}{D'}
 =\frac{C-n'}D+\frac{n'}{D'}\frac{D'-D}D
 \le\frac{4B}{D}.
\]
For $J>0$, put $p=C/J$. Since $p-C/(J+2)=2p/(J+2)$,
\[
 0\le p-\widehat q\le\min\{p,A/(J+2)\}.
\]
The map $x\mapsto\min\{x,c\}$ is $1$-Lipschitz on nonnegative reals. A triangle inequality therefore proves \eqref{bin:pathloss}. If $J=0$, then $C=n'=\widehat q=0$, and $A\ge2$ makes the second term exactly $q_{Y_M}$. Thus zero rows and attack-created false exits are covered without any purity assumption transferred from a different observation model.
\end{proof}

\subsection{The inherited occupation interfaces and their domains}
We explicitly retain three predecessor inputs. These statements are analytic interfaces from the complete predecessor dossier; their earlier proofs are not replaced by a claim that terminal-selected row observations are independent.

First, the original clean add-one theorem R011 gives, for every binary source, every initial pair law, every $\beta\in[0,1]$, and $M\ge64$,
\begin{equation}
 S_{\rm add}\le8192\sqrt{32}/\sqrt{M+1}.
 \label{bin:cleanimport}
\end{equation}
Second, the universal terminal-occupation inequality for a four-state chain is
\begin{equation}
 \mathbb P(J<w)\le\min\{1,16w/(M+1)\},
 \qquad w\in\{1,2,\ldots\},
 \label{bin:universal}
\end{equation}
with arbitrary initial law and no mixing requirement. It implies
\begin{align}
 \mathbb E\frac1{J+2}&\le\frac{\sqrt{32}}{\sqrt{M+1}},\label{bin:reciprocal}\\
 \mathbb E\min\{1,A/(J+2)\}&\le\Psi\left(\frac{16A}{M+1}\right),
 \quad A>0,\label{bin:layer}
\end{align}
where $\Psi(x)=x[1+\log(1/x)]$ for $0<x<1$, $\Psi(x)=1$ for $x\ge1$, and $\Psi(0)=0$. Indeed, at layer height $v>0$, a nonempty event $J<A/v-2$ has integer threshold $\lceil A/v-2\rceil\le A/v$; integrating \eqref{bin:universal} proves \eqref{bin:layer}. The bound $\Psi(x)\le\sqrt{2x}$ gives \eqref{bin:reciprocal}. Since
\[
 \left|\frac CJ-\frac{C+1}{J+2}\right|
 =\frac{|2C/J-1|}{J+2}\le\frac1{J+2}\quad(J>0),
\]
we obtain
\begin{equation}
 \mathbb E\left[\mathbf 1\{J>0\}|C/J-q_{Y_M}|\right]
 \le(8192+1)\sqrt{32}/\sqrt{M+1}.
 \label{bin:comparator}
\end{equation}

Third, the weighted occupation interfaces imported from BIF001, equation (3.4), and BAI003 version 2, equation (6.2), are precisely the following. For the fixed true weights $q=(a,b,e,t)$ and every $A>0$,
\begin{align}
 \mathbb E\min\{q_{Y_M},A/(J+2)\}
 &\le4M^{-1/2}+\frac{16672A}{M}
       \left[1+\log\left(1+\frac1{\sqrt\beta A}\right)\right]
 \label{bin:fullsupport}
\end{align}
when $M\ge64$, $\beta M^2\ge1$, the stationary pair law has full support, and the all-state certificate holds. This is uniform over \emph{every} initial law. Its crossing weights are $b,e$, not unit caps. Its proof uses first diagonal hitting and weighted terminal occupations; it is not a global square-root-$\beta$ mixing estimate. The diagonal-start component bounds used in that proof hold at every suffix horizon $m\ge8$, even if $\beta m^2<1$.

If $00$ is absorbing, with arbitrary remaining row probabilities and arbitrary initial law, then for every $M\ge8$ and $A>0$,
\begin{equation}
 \mathbb E\min\{q_{Y_M},A/(J+2)\}
 \le\frac{256\max\{A,1\}}M+\frac{16}{M+2},
 \qquad q=(0,b,e,t).
 \label{bin:absorbing}
\end{equation}
There is no $\beta$ restriction in this interface. Symbol complementation covers absorbing $11$; it sends $q$ to $(t,e,b,a)$ and preserves the same fixed orientation. In particular, this is not a source-dependent branch of the algorithm. These are the occupation statements being imported, not their predecessor quantized-bank estimators or architecture-restricted lower bounds.

\begin{proposition}[Binary robust upper bound]
\label{bin:upperprop}
The estimator above satisfies, uniformly over all sources, initial laws, and admissible attacks,
\[
 \mathbb E\,\operatorname{TV}(\widehat\theta,\theta_{Y_M})
 \le\min\{1,2^{17}s+2^{19}F_2\}.
\]
\end{proposition}
\begin{proof}
Suppose first $k\ge1$, so $k\le A=24k+2\le26k$. For full support and $\beta\ge M^{-2}$, \eqref{bin:fullsupport} gives
\[
 \mathbb E\min\{q_{Y_M},A/(J+2)\}\le4s+433472F_2.
\]
For positive $\beta$ the non-full-support case is exhaustive as follows. Stationarity has the form $\Pi=(y,x,x,w)$ with $ya=xe$ and $wt=xb$. If $x>0$, the all-state inequalities
\[
 P^2(10,01)=ea\ge\beta x,\qquad P^2(01,10)=bt\ge\beta x
\]
force all four row parameters positive and hence $y,w>0$. If $x=0$, every supported diagonal is absorbing. Two such diagonals cannot both carry stationary mass under a positive all-state certificate. Thus non-full support is a singleton absorbing diagonal. Equation \eqref{bin:absorbing} then bounds the functional by $6672k/M\le6672F_2$, including starts at unsupported or transient states.

Finally, if $0\le\beta<M^{-2}$, no support classification is needed. Now $e_*=1/M$, and \eqref{bin:layer} gives
\[
 \mathbb E\min\{q_{Y_M},A/(J+2)\}
 \le\Psi(416k/M)\le416\frac{k}{M}[1+\log(1+M/k)].
\]
For the last inequality set $z=M/k\ge1$. At $z<416$ the right side exceeds one; at $z\ge416$ substitute the formula for $\Psi(416/z)$ directly. Equality $\beta=M^{-2}$ belongs to the full-support or absorbing strong case. Combining these bounds with \eqref{bin:pathloss} and \eqref{bin:comparator} proves the claim since
\[
 2(8192+1)\sqrt{32}+4<2^{17},\qquad433472<2^{19}.
\]
At $k=0$, the clean branch and \eqref{bin:cleanimport} suffice; no singular logarithm is evaluated. Total variation is always at most one.
\end{proof}

The estimator is a finite rational count algorithm. It can read and retain the corrupted word, identify its final pair, scan the $M$ triples, and maintain two row counters. The output can be represented by a shared positive denominator and two integer numerators; no division, logarithm, or square root is required. All these integers have $O(\log(M+2))$ bits. This observation is not a neural-depth theorem, an attention-budget theorem, or an unknown-$k$ adaptation result.

\section{Binary renewal words and the obstruction to naive erasure}
\label{bin:obstruction}
The lower bound requires accounting for the entire arrangement of the corrupted word. In the binary case, simply deleting rare exits does not erase the timing information carried by variable cycle lengths. The following exact obstruction explains why a different transport argument is needed.

Let $S=01$ and $L=101$. Denote by $U_{r,\ell}$ the uniform distribution on words obtained by concatenating $r$ copies of $S$ and $\ell$ copies of $L$ in every possible order. Their common length is $t=2r+3\ell$. The code is prefix-free, so cycle order is uniquely recoverable, and every nonempty returned word ends in $01$.

\begin{proposition}[An exact boundary obstruction and a biased cheap merge]
\label{bin:obstructionprop}
For $t=2r+6$, the optimal expected Hamming cost of transport from $U_{r,2}$ to $U_{r+3,0}$ is $t/3+1$. Moreover, the rule that uniformly chooses a consecutive pair of long cycles and replaces $L S^g L$ by $S^{g+3}$ has expected cost $3+2r/(\ell+1)$, but does not in general have marginal $U_{r+3,\ell-2}$.
\end{proposition}
\begin{proof}
For $\ell=2$, write a uniformly sampled word as $S^{g_0}LS^{g_1}LS^{g_2}$. The vector $(g_0,g_1,g_2)$ is uniform over weak compositions of $r$. The target is the single word $S^{r+3}$. Outside $LS^{g_1}L$ the two words agree, and direct comparison inside it gives exactly $2g_1+3$ disagreements. Since $\mathbb E g_1=r/3$, every coupling has cost $2r/3+3=t/3+1$.

For general $\ell$, the gap vector $(g_0,\ldots,g_\ell)$ is again a uniform weak composition of $r$. Choosing uniformly among the $\ell-1$ adjacent pairs in the ordered list of long cycles selects an intervening gap with mean $r/(\ell+1)$, proving the stated mean cost. If the output gap vector is $h=(h_0,\ldots,h_{\ell-2})$, the number of preimages obtained by merging at coordinate $i$ is $\binom{h_i-1}{2}$, interpreted as zero for $h_i<3$. Indeed, the three original gaps involved in that merge sum to $h_i-3$. Thus the output density relative to the uniform target is proportional to
\[
 D(h)=\sum_i\binom{h_i-1}{2},
\]
which is not constant. For instance, at $(r,\ell)=(1,3)$, the possible target gap vectors have total four, while $D(4,0)=3$ and $D(2,2)=0$. Cheap average cost alone therefore does not establish parameter-independent erasure.
\end{proof}

The obstruction is a boundary statement, not a disproof of interior constant-step transport. The stronger premise
\begin{equation}
 \begin{gathered}
 \text{for each fixed }\eta>0\text{ there is }C_\eta<\infty
 \text{ such that, if }r,\ell\ge\eta t,\ \ell\ge2,\\
 W_1(U_{r,\ell},U_{r+3,\ell-2})\le C_\eta,
 \qquad t=2r+3\ell,
 \end{gathered}
 \label{bin:unprovedT}
\end{equation}
remains unproved here. The metric is Hamming distance. The following reduction is explicitly conditional on this premise; the subsequent entropy proof does not use it.

\paragraph{Why the interior premise would suffice.}
First, there are absolute $\eta_0>0$ and $t_0$ such that, uniformly for $p\in[1/4,3/4]$ and $t\ge t_0$, the reference chain introduced in the next section has
\begin{equation}
 \mathbb P_p\{Y_t=A,\ r,\ell\ge\eta_0t\}\ge1/32.
 \label{bin:conditionalinterior}
\end{equation}
Indeed, its fifth-step $A$-column is $(2p(1-p),(1-p)^2,p)$, so terminal-$A$ probability is at least $1/16$ for every $t\ge5$. At least $n=\lfloor t/3\rfloor$ complete cycles occur by time $t$, since every cycle has length at most three. Among the first $n$ IID cycle choices, both types occur at least $n/8$ except on an event of probability at most $2\exp(-n/32)$ by Hoeffding's inequality. Choose $t_0$ so that this probability is at most $1/32$ and $n/8\ge t/32$. Subtracting this event from the terminal probability proves \eqref{bin:conditionalinterior} with $\eta_0=1/32$. No IID assertion after conditioning on the fixed-time endpoint is used.

Assume \eqref{bin:unprovedT} for $\eta=\eta_0/2$ and enlarge $C_\eta$ to at least one. Put $b=\lfloor k/(128C_\eta)\rfloor$ and first suppose $b\ge1$. For a returned word of length $t$, set $r_t=t\bmod2$; necessarily $\ell\equiv r_t\pmod2$. Quantize
\[
 h=(\ell-r_t)/2,\qquad h_0=b\lfloor h/b\rfloor,
 \qquad\ell_0=r_t+2h_0.
\]
Starting with counts satisfying $r,\ell\ge2\eta t$, compose at most $b-1$ couplings in \eqref{bin:unprovedT}. Each step reduces the long-cycle count by two and increases the short-cycle count by three. If $t\ge2b/\eta$, all intermediate counts remain at least $\eta t$. Coupling composition and the Hamming triangle inequality bound the expected distance from the original word to the final word by $C_\eta b$. The final marginal is exactly $U_{r_0,\ell_0}$, where $2r_0+3\ell_0=t$.

Define the attack on recognizable reference suffixes with these counts and lengths by sampling this composition and accepting its final output only if its distance from the \emph{original} suffix is at most $k$; use identity elsewhere. This is one kernel on all input words, uses no $p$, and preserves the exact terminal pair. The strict-budget truncation loses probability at most $C_\eta b/k\le1/128$ on each eligible class. Conditional on eligibility and $(t,\ell_0)$, the untruncated output has the common law $U_{r_0,\ell_0}$ even after mixing original counts within its cell. At a fixed $t$ there are at most $1+t/(2b)$ such labels; append one bad label for ineligible inputs. This is a finite-label factorization, not a claim that only that many words are possible.

For full justification of the randomization, sample the original counts and eligibility under their $p$-dependent distribution; generate the ideal output from the common target law with an independent seed; then reverse-sample the original arrangement using the selected coupling conditional on its counts and ideal output. The reverse conditional law is independent of $p$. Thus the prior is independent of the auxiliary seed before label selection. The success mass is at least $1/32-1/128=3/128$ by \eqref{bin:conditionalinterior}. Fixing that seed and decoder randomness leaves at most $2+t/(2b)$ candidate scalar predictions. The union of their sufficiently small radius-$c/(2+t/(2b))$ neighborhoods has uniform-prior mass less than half the success mass. The small-ball argument proved in detail in Lemma~\ref{bin:smallball} therefore gives a reference loss at least $c_1/(1+t/b)$ with an absolute $c_1>0$. It subtracts bad mass from a probability, not directly from a vanishing risk bound.

Use the stationary source and path domination of Section~\ref{bin:witness}, with the same uniform prior and the same attack for every suffix length. Write $H=\lfloor1/(4e_*)\rfloor-1$ and $t_1=\max\{t_0,\lceil2b/\eta\rceil\}$. Since $b\ge1$, there is a constant $C_1$ depending only on $\eta,t_0$ such that $t_1\le C_1b$. The disjoint suffixes yield a lower bound proportional to
\[
 \frac1M\sum_{t=t_1}^{H}\frac1{1+t/b}.
\]
The sum is zero if its range is empty. Otherwise integral comparison bounds it below by
\[
 \frac bM\left[\log(1+H/b)-\log(1+C_1)\right]_+.
\]
For $e_*\le1/16$, $H\ge1/(8e_*)$. Also $b$ and $k$ differ by fixed multiplicative constants when $b\ge1$, because $k/(256C_\eta)\le b\le k/(128C_\eta)$. Consequently the separate IID floor $ck/M$ absorbs the omitted initial range and the bounded difference of logarithms, giving the complete corruption profile up to constants depending only on $C_\eta,\eta,t_0$. If $b=0$, then $k<128C_\eta$ is bounded: using $e_*\ge1/M$,
\[
 \frac{k}{\sqrt M}[1+\log(1+M/k)]
 \le k\,[1+\log2+2/\exp(1)],
\]
so a constant multiple of the clean floor covers the full corruption term. For $\beta>1/256$, the inequality $F_2\le k/M+2M^{-1/2}$ and the two IID floors suffice directly. This completes only the implication from \eqref{bin:unprovedT}. The stronger premise remains open; the proved route below replaces it by a weaker entropy transport estimate with a genuine length cutoff.

\section{A stationary binary witness with an all-state certificate}
\label{bin:witness}
For $\beta\le1/256$ set
\[
 a=\frac1{2M},\qquad e_*=\max\{\sqrt\beta,1/M\},\qquad
 p\in[1/4,3/4].
\]
With pair states $W=00,A=01,B=10,C=11$, define the source by
\begin{equation}
 \bigl(\theta_W(1),\theta_A(1),\theta_B(1),\theta_C(1)\bigr)
       =(a,p,1-e_*,e_*).
 \label{bin:source}
\end{equation}
This is a modified witness, not an identification with a predecessor binary precision witness. Put
\begin{equation}
 x=\left(e_*/a+2+\frac p{1-e_*}\right)^{-1},\qquad
 \Pi=x\left(e_*/a,1,1,\frac p{1-e_*}\right).
 \label{bin:stationary}
\end{equation}

\begin{lemma}[Stationarity and minorization]
\label{bin:certificate}
The law \eqref{bin:stationary} is stationary and has $\Pi_W\ge1/4$. Its pair kernel obeys $P^2(u,v)\ge\beta\Pi(v)$ for every pair of states $u,v$.
\end{lemma}
\begin{proof}
Direct balance gives the stationary equations: the crossing masses are equal, $\Pi_W a=\Pi_B e_*$, and $\Pi_C(1-e_*)=\Pi_Ap$. Normalization gives \eqref{bin:stationary}. Since $e_*/a\ge2$, $e_*\le1/16$, and $2+p/(1-e_*)\le2.8$, we have $\Pi_W\ge2/4.8>1/4$ and $x\le a/e_*$. The two-step matrix in the stated order is
\[
 P^2=\begin{pmatrix}
 (1-a)^2&(1-a)a&a(1-p)&ap\\
 (1-p)e_*&(1-p)(1-e_*)&p(1-e_*)&pe_*\\
 e_*(1-a)&e_*a&(1-e_*)(1-p)&(1-e_*)p\\
 e_*(1-e_*)&(1-e_*)^2&e_*(1-e_*)&e_*^2
 \end{pmatrix}.
\]
Here $a\le1/128$, $\beta\le e_*^2$, and $\beta\le e_*/16$. For column $W$, every entry is at least $\beta$, using $(1-p)e_*\ge e_*/4$ and $e_*(1-e_*)\ge15e_*/16$, hence dominates $\beta\Pi_W$. In column $A$, the small entry $e_*a$ dominates $\beta a/e_*\ge\beta x$, while $(1-a)a\ge\beta a/e_*$ and the other entries are bounded below by fixed positive constants. For column $B$, the first-row entry is at least $a/4\ge\beta a/e_*$, the last-row entry is at least $\beta$, and the remaining entries have constant lower bounds. For column $C$, the first-row comparison is
\[
 ap\ge\frac{\beta pa}{e_*(1-e_*)}\ge\frac{\beta px}{1-e_*}=\beta\Pi_C,
\]
using $e_*(1-e_*)\ge\beta$. The second and third rows satisfy the same target comparison because $e_*,1-e_*\ge\beta x/(1-e_*)$; the last row uses $e_*^2\ge\beta\ge\beta\Pi_C$. This verifies the certificate even on states that an initialization might not visit.
\end{proof}

The active reference chain $R_p$ on $A,B,C$ has
\[
 R_p(A,B)=1-p,\quad R_p(A,C)=p,\quad R_p(B,A)=R_p(C,B)=1.
\]
Its return cycles to $A$ emit precisely $S$ and $L$. If $\nu_{p,t}$ denotes its law of $t$ emitted bits from $A$, the cycle-count-conditioned law is $U_{r,\ell}$, since every ordering has probability $(1-p)^r p^\ell$.

Initialize the actual source at $\Pi$. Let $E_j$ consist of a prefix of $N-j$ zeros, followed by the first $1$, and then $j-1$ reference transitions. These events, over the different $j$, are disjoint. Relative to the reference law of the suffix, each actual word has probability multiplier
\begin{equation}
 \Pi_W(1-a)^{M-j}a(1-e_*)^{N_B+N_C},
 \label{bin:density}
\end{equation}
where $N_B,N_C$ count reference departures. For $j\le1/(4e_*)$, this multiplier is at least $1/(32M)$: use $\Pi_W\ge1/4$, $(1-a)^M\ge1/2$, $a=1/(2M)$, and $(1-e_*)^j\ge1-je_*\ge3/4$. This is pointwise path domination. The probabilities of $E_j$ can depend on $p$; we will not declare a prior uniform after conditioning on $E_j$. At terminal $A$ the true actual prediction target is Bernoulli$(p)$.

\section{Entropy transport to a public canonical bridge}
\label{bin:transport}
We now prove the weaker transport estimate that is sufficient for the full minimax rate. For probability measures on a finite bit cube, $W_1$ denotes minimum expected Hamming distance over all couplings.

\begin{lemma}[Uniform coupling, concentration, and entropy transport]
\label{bin:entropy}
Let $p,q\in[1/8,7/8]$ and $t\ge1$. With $L_0=326$, every Hamming-$1$-Lipschitz function $f$ satisfies
\begin{equation}
 \log\mathbb E_{\nu_{p,t}}
       \exp\{\lambda(f-\mathbb E_{\nu_{p,t}}f)\}
       \le tL_0^2\lambda^2/8.
 \label{bin:mgf}
\end{equation}
Consequently, for $\mu\ll\nu_{p,t}$,
\begin{align}
 W_1(\mu,\nu_{p,t})&\le L_0\sqrt{t\,\operatorname{KL}(\mu\Vert\nu_{p,t})/2},
 \label{bin:transportentropy}\\
 W_1(\nu_{p,t},\nu_{q,t})&\le L_0t|p-q|.
 \label{bin:transportparameter}
\end{align}
\end{lemma}
\begin{proof}
The $A$-column of $R_p^5$ is
\[
 \bigl(2p(1-p),(1-p)^2,p\bigr),
\]
whose entries are at least $1/64$. Two chains with the same parameter and arbitrary initial states can therefore couple at five-step boundaries using this common atom. Conditional bridges fill in each block with the correct marginal law; after meeting, use the same transitions. The expected number of emitted-bit disagreements, also for a truncated final block, is at most
\[
 5\sum_{r\ge0}(63/64)^r=320.
\]
Reveal the transitions one at a time in the Doob martingale of $f$. Two possible current transitions change at most one current bit, and their conditional tails can be coupled with expected cost at most $320$. Thus the conditional range of an increment is at most $321\le L_0$. Conditional Hoeffding's lemma followed by iteration gives \eqref{bin:mgf}. If $f$ is initially defined only on the path support, a Hamming-Lipschitz extension to the cube supplies the same conclusion.

For $\lambda>0$, the finite entropy inequality is
\[
 \mathbb E_\mu f-\mathbb E_\nu f
 \le\frac{\operatorname{KL}(\mu\Vert\nu)}\lambda
       +\frac{tL_0^2\lambda}{8},\qquad\nu=\nu_{p,t}.
\]
To check it directly, set
\[
 Q(x)=\frac{\nu(x)e^{\lambda f(x)}}{\mathbb E_\nu e^{\lambda f}}.
\]
Use nonnegativity of $\operatorname{KL}(\mu\Vert Q)$ and then \eqref{bin:mgf}. Optimizing in $\lambda$ gives $L_0\sqrt{t\operatorname{KL}/2}$, including the zero-entropy case by a limit. Finite Kantorovich--Rubinstein linear-programming duality identifies $W_1$ with the supremum of $\mathbb E_\mu f-\mathbb E_\nu f$ over Hamming-$1$-Lipschitz $f$, proving \eqref{bin:transportentropy}. The finite metric formulation permits zero-probability ambient words through the same extension.

Finally, interpolate from the $p$ law to the $q$ law by changing one transition kernel at a time, with a common preceding segment and a common-$q$ tail. The changed transition disagrees with probability at most $|p-q|$. Conditional on that disagreement, the current bit and coupled continuation cost at most $321\le L_0$. Summing the $t$ hybrid comparisons proves \eqref{bin:transportparameter}.
\end{proof}

\begin{lemma}[Microcanonical-to-canonical bridge bound]
\label{bin:bridge}
Suppose $t=2r+3\ell\ge5$, $m=r+\ell$, and $\widehat p=\ell/m\in[1/8,7/8]$. For any $q\in[1/8,7/8]$, let $\nu^A_{q,t}$ be $\nu_{q,t}$ conditioned on terminal pair $A$. Then
\begin{equation}
 W_1(U_{r,\ell},\nu^A_{q,t})
 \le1024\sqrt{t\log(t+1)}+326t|\widehat p-q|.
 \label{bin:bridgebound}
\end{equation}
Both endpoint measures in this transport have final two bits $01$.
\end{lemma}
\begin{proof}
The event of exactly these cycle counts and a return at time $t$ has $\nu_{\widehat p,t}$-probability
\[
 \binom m\ell\widehat p^\ell(1-\widehat p)^r\ge\frac1{m+1}\ge\frac1{t+1}.
\]
The equality follows from prefix-free decoding; the first inequality holds because $\ell$ is a mode of Binomial$(m,\ell/m)$ and the distribution has $m+1$ possible values. Conditioning on that event is exactly $U_{r,\ell}$, whence
\[
 \operatorname{KL}(U_{r,\ell}\Vert\nu_{\widehat p,t})\le\log(t+1).
\]
Likewise $R_q^5(u,A)\ge1/64$ gives terminal-$A$ probability at least $1/64$ for every $t\ge5$, so
\[
 \operatorname{KL}(\nu^A_{q,t}\Vert\nu_{q,t})\le\log64.
\]
Apply \eqref{bin:transportentropy} to these two conditionings, \eqref{bin:transportparameter} to the intervening parameters, and the triangle inequality and symmetry of $W_1$. The result is
\[
 L_0\sqrt{t/2}\bigl(\sqrt{\log(t+1)}+\sqrt{\log64}\bigr)
      +L_0t|\widehat p-q|.
\]
Since $326(1+\sqrt6)/\sqrt2<1024$, this is bounded by \eqref{bin:bridgebound}. Unconditioned intermediate measures are used only in this metric triangle inequality; they are not possible final corrupted outputs.
\end{proof}

\section{A single strict-budget binary attack and its information bound}
\label{bin:attack}
Fix $k\ge1$ and set $c_E=2^{-44}$. For each $t\ge4096$, partition $[1/8,7/8]$ into
\[
 n_t=\max\{1,\lceil3\cdot2^{18}t/k\rceil\}
\]
equal subintervals. The resulting rational grid has spacing at most $\Delta_t=k/(2^{20}t)$ and at most $2+2^{20}t/k$ points. Round $\widehat p$ to the nearest grid point $q$, breaking ties deterministically. For
\begin{equation}
 4096\le t\le c_E k^2/\log(k+2),
 \label{bin:entropycutoff}
\end{equation}
we have $t\le k^2$ and $\log(t+1)\le2\log(k+2)$. Hence \eqref{bin:bridgebound} is at most
\begin{equation}
 \left(1024\sqrt{2c_E}+\frac{326}{2^{20}}\right)k
       <k/1024.
 \label{bin:costsmall}
\end{equation}

For every eligible public count pair $(r,\ell)$, choose a minimum-cost coupling between $U_{r,\ell}$ and $\nu^A_{q,t}$. This is a finite transport problem with rational marginal probabilities; a rational optimal vertex exists. A fixed lexicographic rule can select one. The table depends on $t,k,r,\ell,q$, never on the unknown parameter $p$.

The attack is defined on every binary word of length $N=M+2$. If the word has at least two initial zeros, then a first $1$, and thereafter follows the reference transitions and ends in pair $A$, let $t$ be the number of transitions after that first $1$. Uniquely decode its cycle counts. If \eqref{bin:entropycutoff} holds and $\widehat p\in[1/8,7/8]$, sample a candidate $Y$ through the chosen coupling conditional on the observed suffix $X$. Replace the suffix by $Y$ only when $d_H(X,Y)\le k$; otherwise leave the entire word unchanged. On all other input words, also use identity corruption. This is one source-independent offline randomized kernel. It preserves the zero prefix, the first $1$, and the exact final pair. All changed positions are among the first $M$, and the budget holds on every realization, not only in expectation. By \eqref{bin:costsmall} and Markov's inequality, truncation failure on any eligible microcanonical input class is at most $1/1024$.

\begin{lemma}[Independent-seed small-ball lower bound]
\label{bin:smallball}
Let $p$ have the single uniform prior on $I=[1/4,3/4]$. For every fixed $t$ satisfying \eqref{bin:entropycutoff}, every possibly randomized predictor of the complete corrupted word has reference loss, restricted to successful terminal-$A$ paths, at least
\begin{equation}
 \frac{2^{-38}}{1+t/k}.
 \label{bin:referencelower}
\end{equation}
\end{lemma}
\begin{proof}
For $p\in I$, the $A$-column of $R_p^5$ is bounded below by $1/16$, so the reference terminal-$A$ probability is at least $1/16$ for $t\ge5$. On a return at time $t$, the number $m$ of complete cycles lies between $\lceil t/3\rceil$ and $\lfloor t/2\rfloor$. Before imposing the return event, the successive cycle choices are IID Bernoulli$(p)$. For each fixed possible $m$, Hoeffding's inequality bounds the chance that its first $m$ choices have empirical fraction outside $[1/8,7/8]$ by $2e^{-m/32}$. A union over at most $t$ deterministic values of $m$ gives a bound $2t e^{-t/96}\le1/128$ for $t\ge4096$; the expression is decreasing on this range. Thus eligible endpoint-$A$ mass is at least $1/32$, uniformly in $p$. This argument does not condition the cycle choices into an IID sample at a stopping time. Removing truncation failures leaves success probability at least $1/32-1/1024\ge1/64$.

Consider an auxiliary untruncated observation. On eligible count classes use the sampled bridge $Y$ without the budget test. On all other paths output a fixed public length-$t$ returned word ending in $01$. Such a word exists for every $t\ge2$, since every integer $t\ge2$ is a nonnegative combination of two and three. Prepend the public prefix appropriate to this fixed $t$. The auxiliary observation factors through the selected grid point, or one additional bad label, followed by a channel independent of $p$. Its number of labels is at most
\begin{equation}
 D_t=3+2^{20}t/k\le2^{22}(1+t/k).
 \label{bin:labels}
\end{equation}
It is the label count, not the number of possible output words, that is finite here.

For precision about the randomization, first draw the original counts and eligibility under their $p$-dependent law. Then generate the bridge for the selected label using a uniform seed $U$ independent of $p$. Finally, reverse-sample the original suffix from the selected coupling conditional on its counts and this bridge. That reverse disintegration depends on the counts and the coupling, but not on $p$, because the microcanonical law is uniform. The bad label uses the fixed word. This construction reproduces the original joint experiment and its budget-success event. For ineligible counts, the original conditional path can be sampled separately; no agreement is claimed on that branch.

Fix $U$ and the predictor's independent randomization. There are now at most $D_t$ candidate predicted probabilities of symbol $1$, one for each possible auxiliary label. With
\[
 \rho=\frac1{512D_t},
\]
the union of their radius-$\rho$ neighborhoods has uniform-prior mass at most $2D_t\rho/|I|=1/128$. This bound does not depend on how a parameter-dependent stochastic encoder chooses its label. The prior remains independent of $U$ before that selection. Subtract this neighborhood probability from success mass at least $1/64$. With probability at least $1/128$, the actual and auxiliary observations agree, their true endpoint is $A$, and the prediction error is at least $\rho$. Averaging the fixed seeds yields loss at least
\[
 \frac1{65536D_t}\ge\frac{2^{-38}}{1+t/k}.
\]
The constant failure probability has been subtracted from a probability bound before multiplying by $\rho$; subtracting it directly from an order-$D_t^{-1}$ risk bound would be invalid. The target on the retained event is the true Bernoulli$(p)$ row, and no posterior after terminal or count conditioning has been called uniform.
\end{proof}

For the actual stationary witness, put
\begin{equation}
 J_* =\left\lfloor\min\left\{\frac1{8e_*},
                      \frac{2^{-44}k^2}{\log(k+2)}\right\}\right\rfloor,
 \qquad
 T=\frac1M\sum_{t=4096}^{J_*}\frac1{1+t/k},
 \label{bin:sumdefinition}
\end{equation}
with an empty sum defined as zero. If $4096\le t\le J_*$ and $j=t+1$, then $j\le1/(4e_*)$ because $e_*\le1/16$, and $j\le M/8+1\le M/4$. The recognizable zero prefix is therefore present, and the density domination \eqref{bin:density} applies. Use the same prior, actual source, stationary initialization, and attack kernel for every disjoint $E_j$. Summing \eqref{bin:referencelower} and discarding all other nonnegative losses gives
\begin{equation}
 R_{\rm stat}\ge2^{-43}T.
 \label{bin:sumbound}
\end{equation}
The source supremum is not moved inside a suffix sum. Pointwise domination justifies the common-prior integration even though the rare-entry event probabilities depend on the source parameter.

\section{Exact scalar completion and the binary minimax lower bound}
\label{bin:completion}
The transport cutoff is genuine. It is removed from the final profile by a finite scalar inequality, rather than by extending the transport argument to longer suffixes.

\begin{lemma}[Completion with floors and empty sums]
\label{bin:scalar}
Let $M\ge64$, $1\le k\le M$, $e_*\ge1/M$, $c>0$, $C\ge1$, and let $j_0$ be a positive integer. Set
\[
 \epsilon=k/M,\quad s=M^{-1/2},\quad
 U=\min\{1/(Ce_*),ck^2/\log(k+2)\},\quad
 T_0=M^{-1}\sum_{j=j_0}^{\lfloor U\rfloor}(1+j/k)^{-1},
\]
where an empty sum is zero. Then
\begin{equation}
 \epsilon\left[1+\log\left(1+\frac1{e_*k}\right)\right]
 \le D(c,C)[1+\log(1+j_0)]\max\{s,\epsilon,T_0\},
 \label{bin:scalarbound}
\end{equation}
where
\[
 D(c,C)=7+\log\max\{1,C\}+6\log\max\{1,2/c\}.
\]
In particular, for $c=2^{-44}$, $C=8$, and $j_0=4096$, the coefficient in \eqref{bin:scalarbound} is less than $1830$.
\end{lemma}
\begin{proof}
Put $A_0=1/(e_*k)$, $B_0=ck/\log(k+2)$, and
$L=\log(1+\min\{A_0/C,B_0\})$. We first prove
\begin{equation}
 \epsilon[1+\log(1+A_0)]
 \le D(c,C)\max\{s,\epsilon,\epsilon L\}.
 \label{bin:scalarintermediate}
\end{equation}
If $k\le M^{1/4}$, use $A_0\le M/k\le M$ to obtain
\[
 \frac{\epsilon[1+\log(1+A_0)]}{s}
 \le M^{-1/4}[1+\log(1+M)]
 \le1+\log2+4/\exp(1)<4.
\]
Here the maximum of $x^{-1/4}\log x$ on $x\ge1$ is $4/\exp(1)$, and $D(c,C)\ge7$.

If $k>M^{1/4}$, then $A_0\le M/k<k^3$. Also
\[
 \log(k+2)\le2\sqrt k\quad(k\ge1),
\]
because the difference is positive at one and its derivative $1/\sqrt k-1/(k+2)$ is positive. Let $q_0=\max\{1,2/c\}$. It follows that $1+\sqrt k\le q_0(1+B_0)$, and hence
\[
 \log(1+A_0)\le3\log(1+k)
 \le6\log(1+\sqrt k)
 \le6\log q_0+6\log(1+B_0).
\]
If $A_0/C\le B_0$, then $1+A_0\le\max\{1,C\}(1+A_0/C)$, so the left side of \eqref{bin:scalarintermediate} is at most
\[
 [2+\log\max\{1,C\}]\max\{\epsilon,\epsilon L\}.
\]
If $A_0/C>B_0$, the preceding logarithmic bound instead gives
$[7+6\log q_0]\max\{\epsilon,\epsilon L\}$. Both coefficients are at most $D(c,C)$, proving \eqref{bin:scalarintermediate} throughout the finite domain.

Now $L=\log(1+U/k)$. If $J=\lfloor U\rfloor\ge j_0$, monotonicity gives
\[
 \frac{T_0}{\epsilon}
 =\sum_{j=j_0}^{J}\frac1{k+j}
 \ge\int_{j_0}^{J+1}\frac{dx}{k+x}
 \ge L-\log(1+j_0/k).
\]
If $J<j_0$, then $U<j_0$ and directly $L\le\log(1+j_0/k)$. Thus in both cases
\[
 \epsilon L\le T_0+\epsilon\log(1+j_0)
 \le[1+\log(1+j_0)]\max\{\epsilon,T_0\}.
\]
Combining with \eqref{bin:scalarintermediate} proves \eqref{bin:scalarbound}. At the specified constants, its coefficient is
\[
 [7+273\log2][1+\log4097]<1830.
\]
No requirement that the sum be nonempty, that $k$ diverge, or that the logarithm be large was used.
\end{proof}

For completeness, the separate lower floors used in the completion are full-history, stationary IID witnesses. They have all four rows equal to Bernoulli$(p)$ and stationary product pair law, so $P^2$ has every row exactly $\Pi$ and satisfies every certificate $\beta\in[0,1]$.

For the clean floor take identity corruption and
\[
 p_\pm=1/2\pm\frac1{16\sqrt{N}},\qquad N=M+2.
\]
The elementary Bernoulli bound $\operatorname{KL}(p\Vert q)\le(p-q)^2/[q(1-q)]$ gives full-word KL at most $5/64$, and Pinsker gives total variation below $1/4$. Nearest-target testing, or the triangle inequality integrated against the common part of the two observation laws, implies the conservative bound
\begin{equation}
 R_{\rm stat}\ge\frac1{256\sqrt M}.
 \label{bin:cleanfloor}
\end{equation}
The identity attack is admissible for every public budget $k$.

For the corruption floor take the uniform prior on $[1/4,3/4]$ and one common kernel. If there are $n$ ones among the first $M$ IID bits, retain
\[
 r=(k+1)\left\lfloor\frac n{k+1}\right\rfloor
\]
ones chosen uniformly among those $n$, changing the rest to zero and preserving the final two bits $V$. The edit count is at most $k$. Conditional on $n$, each specified output arrangement of weight $r$ has probability
\[
 \frac{\binom{M-r}{n-r}}{\binom Mn\binom nr}
       =\frac1{\binom Mr}.
\]
Thus the entire output factors through $(r,V)$, at most
$D=4(\lfloor M/(k+1)\rfloor+1)$ labels, followed by a parameter-independent channel. For a uniform prior of interval length $\ell_0$, any such $D$-message experiment has absolute-error Bayes risk at least $\ell_0/(4D)$. To see this, fix independent channel and decoder randomization; the union of radius-$v$ intervals about the $D$ possible predictions has prior mass at most $2Dv/\ell_0$. Integrate the complementary probability for $0\le v\le\ell_0/(2D)$ and then average the seeds. With $\ell_0=1/2$ and $b=k+1$ this proves
\begin{equation}
 R_{\rm stat}\ge\frac1{32(1+M/b)}\ge\frac{k}{64M},
 \qquad1\le k\le M.
 \label{bin:corruptionfloor}
\end{equation}
The last two bits were included in the label, not conditioned away.

\begin{proposition}[Binary stationary lower bound]
\label{bin:lowerprop}
For every $M\ge64$, $0\le k\le M$, and $\beta\in[0,1]$,
\[
 R_{\rm stat}\ge2^{-54}\max\{M^{-1/2},F_2\}.
\]
\end{proposition}
\begin{proof}
For $\beta\le1/256$ and $k\ge1$, \eqref{bin:sumbound}, \eqref{bin:cleanfloor}, and \eqref{bin:corruptionfloor} are three separate admissible minimax lower bounds. Their maximum gives
\[
 R_{\rm stat}\ge2^{-43}\max\{s,k/M,T\}.
\]
Lemma~\ref{bin:scalar} with the exact sum \eqref{bin:sumdefinition} gives
$\max\{s,F_2\}\le1830\max\{s,k/M,T\}$. Since $1830<2048=2^{11}$, the claimed constant follows. This takes a maximum of different witness bounds, not a sum of their risks.

If $\beta>1/256$, use $\log(1+x)\le x$ and $e_*>1/16$ to obtain
\[
 F_2\le k/M+16/M\le k/M+2s.
\]
The two IID floors then imply $\max\{s,F_2\}\le576R_{\rm stat}$, more than enough. At $k=0$ use only \eqref{bin:cleanfloor} and define $F_2=0$.
\end{proof}

Stationary sources are a subset of sources with arbitrary initial law. Combining this lower bound with Proposition~\ref{bin:upperprop} proves the binary part of the main rate. The attack is not asserted efficient or causal, its seed and edit mask are hidden, and it never edits the terminal pair. The result does not establish the stronger interior transport premise \eqref{bin:unprovedT}. The entropy bridge estimate and its cutoff are compatible with Proposition~\ref{bin:obstructionprop}: they control interior empirical parameters, rather than all neighboring count classes.


\section{Source and version record}\label{app:provenance}
The scientific source set is the stage-eight freeze dated 2 November 2026, 09:22:25.495805 UTC. It comprises SC001 (including its stationary corollary), SC002, SC003, BT001, BT002 (including finite scalar completion), and EP001. Historical candidate headers in the original files are preserved; adoption is established separately by the hash-bound acceptance receipts. In particular, SC003 alone established only the binary upper and elementary floors; the matching logarithmic lower is supplied by BT002. BT001 supplies the obstruction, modified source and a \emph{conditional} reduction, not a proof of its stronger interior constant-cost transport premise.

The new tolerance corollary in Section~\ref{sec:tolerance} is an elementary derivation from the frozen rate formulas and is identified separately for manuscript review. Later exploratory work is not used. The immutable predecessor manuscript, its proofs and the earlier delivery are preserved rather than overwritten. The adjacent \texttt{SOURCE\_MANIFEST.json} records exact file identities, acceptance receipts, inherited interfaces and primary-source reading scope; \texttt{PRESERVATION\_BINDINGS.json} records the pre-authoring preservation baseline. Full predecessor arguments are accessible through these local dossier paths, not through invented publication references. This is an English companion manuscript; no authorship affiliation, external submission, public release or novelty certification is implied.

\bibliographystyle{plainnat}
\bibliography{references}
\end{document}
